% --- Archon physics formalization source begin ---
% archon:physics
% archon:chemistry
% archon:covers IChO2026Problems/problem_icho_2026_t4_a4.lean
% archon:source-report reports/icho_2026/problem_icho_2026_t4_a4.source.json
% archon:problem-id icho_2026_t4
% archon:part-id T4-A4
% archon:previous-part icho_2026_t4_a3
% archon:previous-part-policy natural_language_prerequisite_only

\section{Icho Chemistry problem icho\_2026\_t4\_a4}

\paragraph{Problem source.}
\#\# Source
58th International Chemistry Olympiad, Tashkent, Uzbekistan, 2026.
Official-materials index: https://www.icho2026.uz/problems
English problem PDF: ../icho\_2026\_source/raw/theory\_problem.pdf
English solution/rubric PDF: ../icho\_2026\_source/raw/theory\_solution.pdf
Problem PDF page 38; printed local question page 2; solution starts on PDF page 34.
The page PNG is primary visual evidence for formulas, structures, tables, and figures.

\#\# T4. The Nuclear Past of Uzbekistan
T4. The nuclear past of Uzbekistan 6\% Natural uranium consists primarily of two isotopes, 235U and 238U . 235U is responsible for sustaining nuclear fission reactions. 4.1 Calculate the atomic abundance of 235U in natural uranium, assuming it consists only of isotopes 235U (235.04 a.u.) and 238U (238.05 a.u.). The main reaction occurring in nuclear power plants is the reaction where 235U absorbs a neutron and undergoes fission, releasing energy and producing three additional neutrons, thus sustaining a chain reaction. 235 92 U +1 0 n →… + 31 0n Neutrons that induce the chain reaction are produced by several ways: a) Reaction of 11B with 𝛼-particles inside the reactor core b) Interaction of 𝛾rays with 2D nuclei present in the reactor coolant water c) Reaction of 9Be with 𝛼-particles produced in an externally installed neutron source. 4.2 Write the nuclear reaction equations for (a), (b), and (c). The graph above displays the percentage yield, P(A), of fission products by mass number, A, resulting from 235U fission. 4.3 Write the most common nuclear fission equation for splitting of 235U upon neutron absorption. The pair of elements formed in this reaction are in the same group of Periodic Table.

\#\# Current subquestion 4.4
Calculate the energy released in this reaction (ΔE, MeV), if the binding energy 𝐵𝐸(235U) = 7.59 MeV/nucleon and average binding energy for fission products 𝐵𝐸(fis.) = 8.45 MeV/nucleon. Neglect the binding energy of free neutrons.

\#\# Dataset metadata
IChO 2026; T4; raw rubric points=2; kind=theory; formalization\_ready=true.

\paragraph{Shared problem context.}
T4. The nuclear past of Uzbekistan 6\% Natural uranium consists primarily of two isotopes, 235U and 238U . 235U is responsible for sustaining nuclear fission reactions. 4.1 Calculate the atomic abundance of 235U in natural uranium, assuming it consists only of isotopes 235U (235.04 a.u.) and 238U (238.05 a.u.). The main reaction occurring in nuclear power plants is the reaction where 235U absorbs a neutron and undergoes fission, releasing energy and producing three additional neutrons, thus sustaining a chain reaction. 235 92 U +1 0 n →… + 31 0n Neutrons that induce the chain reaction are produced by several ways: a) Reaction of 11B with 𝛼-particles inside the reactor core b) Interaction of 𝛾rays with 2D nuclei present in the reactor coolant water c) Reaction of 9Be with 𝛼-particles produced in an externally installed neutron source. 4.2 Write the nuclear reaction equations for (a), (b), and (c). The graph above displays the percentage yield, P(A), of fission products by mass number, A, resulting from 235U fission. 4.3 Write the most common nuclear fission equation for splitting of 235U upon neutron absorption. The pair of elements formed in this reaction are in the same group of Periodic Table.

\paragraph{Current subquestion.}
Calculate the energy released in this reaction (ΔE, MeV), if the binding energy 𝐵𝐸(235U) = 7.59 MeV/nucleon and average binding energy for fission products 𝐵𝐸(fis.) = 8.45 MeV/nucleon. Neglect the binding energy of free neutrons.

\paragraph{Recorded answer/context.}
Δ𝐸= 8.45 × 233 −7.59 × 235 = 185.20 MeV Released energy: 185.2 MeV 2 points for the correct answer and relevant calculations; 1 point may be given for the right formula but the wrong answer. 2.0 pt

\paragraph{Figure/image paths.}
\begin{itemize}
\item ../icho\_2026\_source/image/T4\_page-2.png
\item ../icho\_2026\_source/image/T4\_page-1.png
\end{itemize}

\paragraph{Reusable previous-part conclusions.}
\begin{itemize}
\item Source T4-A3. Question: Write the most common nuclear fission equation for splitting of 235U upon neutron absorption. The pair of elements formed in this reaction are in the same group of Periodic Table. Reusable conclusions: Let us presume that total atomic mass of elements in the fission products is ∑𝐴= 235+1−3 = 233 There are two maxima on the graph at 𝐴= 93 and 𝐴= 140, which sum up to 233. Thus, there are only two fission products beside neutrons considered in the question. To find the atomic numbers we use conservation of charge and come to equation: 𝑥+𝑦= 92 with 𝑥and 𝑦denoting atomic numbers of elements. These elements are in the same group, the intervals between them may be 2, 8, 18, 32, and 50. Let’s solve the equation for those parameters: 𝑥+ (𝑥+ 𝑎) = 92 where 𝑎= 2, 8, 18, 32, 50 𝑥= 45, 𝑦= 47 (𝑅ℎ, 𝐴𝑔, 𝑎= 2) 𝑥= 42, 𝑦= 50 (𝑀𝑜, 𝑆𝑛, 𝑎= 8) 𝑥= 37, 𝑦= 55 (𝑅𝑏, 𝐶𝑠, 𝑎= 18) 𝑥= 30, 𝑦= 62 (𝑍𝑛, 𝑆𝑚, 𝑎= 32) 𝑥= 21, 𝑦= 71 (𝑆𝑐, 𝐿𝑢, 𝑎= 50) Technically, Sc-Lu pair might do it, but there is no evidence for such high mass number isotope of scandium (𝐴= 93), so the only reasonable answer is Rb and Cs. Equation of nuclear fission reaction: 235 92 U + 1 0n →93 37Rb + 140 55 Cs + 31 0n 3 points total: 0.6 point for each correct atomic mass of two elements (1.2 points). If the student provides the sum of the atomic masses of two elements they will also be given 1.2 points. 0.6 point for each correct element (1.2 points). 0.6 point for nuclear fission reaction equation. 3.0 pt Policy: natural\_language\_prerequisite\_only; the current target and later answers must not be assumed
\end{itemize}

\paragraph{Formalization target.}
create a compiling Lean file with sorry bodies at `IChO2026Problems/problem\_icho\_2026\_t4\_a4.lean`.
The Lean declarations must preserve every mathematical object, quantifier, hypothesis, side condition, approximation/error guarantee, and requested conclusion from the source statement.
Use Mathlib/Physlib/Chemistry names found through LeanExplore where available. Prefer the configured benchmark Base library; introduce a local definition only when the source genuinely requires an object that the environment does not expose.

\begin{theorem}\leanok
[Icho Chemistry formalization target]
\label{thm:physics:icho_2026_t4_a4:target}
\lean{IChO2026Problems.T4A4.uranium235_fission_released_energy}
The supplied binding energies determine the energy released by the selected
$^{235}$U fission reaction.
\end{theorem}
\begin{proof}

\leanok
Apply nucleon-number conservation, multiply the respective average binding
energies by the source and fission-product nucleon totals, and subtract the
initial total from the final total.
\end{proof}
% --- Archon physics formalization source end ---
